% ============================================================
%  Chapter 6 — Conclusion
% ============================================================
\chapter{Conclusion and Future Work}
\label{ch:conclusion}

\section{Summary of the Study}
\label{sec:concl_summary}

This thesis set out to bring groundwater potential mapping of the \studyarea{}
from a single district and a manual method up to the full five-district plateau
and a machine-learning method. Nine geo-environmental features---elevation,
slope, aspect, topographic wetness index, NDVI, NDWI, rainfall, soil texture, and
land surface temperature---were assembled at 30~m resolution through Google Earth
Engine for the whole of Chakwal, Rawalpindi, Attock, Jhelum, and Mianwali,
yielding \npixels{} valid pixels. A knowledge-based weighted overlay supplied
three-class reference labels, a balanced sample of 450,000 pixels was used to
train four classifiers, and the best of these was applied across the plateau to
produce a full-resolution map, which was then tested against independent evidence.
The work thereby answers the four research questions posed in
Chapter~\ref{ch:introduction}, as summarized below.

\section{Key Findings}
\label{sec:concl_findings}

\begin{itemize}
  \item \textbf{Reproducing and refining the overlay (RQ1).} A single XGBoost
        model reproduced the knowledge-based weighted overlay across all
        37~million pixels with 99.05\% agreement, confirming that supervised
        learning can stand in for the manual method---while doing so in a form
        that needs no fixed expert weights and yields a continuous probability
        surface.

  \item \textbf{Best model (RQ2).} Of the four classifiers, XGBoost performed best
        on every metric (99.12\% accuracy, macro $F_1 = 0.991$, Cohen's
        $\kappa = 0.987$, ROC-AUC $= 0.9998$), narrowly ahead of LightGBM and the
        voting ensemble and a clear step above random forest. The differences
        among the boosting models are small enough that any of them would serve in
        practice.

  \item \textbf{Controlling factors (RQ3).} Soil texture ($r = 0.836$; XGBoost
        importance 0.831) and rainfall ($r = 0.514$) emerged as the dominant
        predictors from both the correlation analysis and the trained models,
        consistent with the hydrogeology of a semi-arid, recharge-limited plateau
        and with the reference study's emphasis on permeability-related factors.

  \item \textbf{Independent credibility (RQ4).} A spatial-realism validation, using
        evidence that played no part in training, supported the map on all four
        checks: known alluvial sites fell in Medium or High ground, independent
        rain-gauge stations separated by hydrogeological setting at a significant
        level ($p = 0.023$), the 20\% High fraction matched independent regional
        AHP work, and the district-scale potential gradient reproduced the expected
        north--south ordering exactly ($\rho = 1.00$).
\end{itemize}

Across the plateau the model classified 20\% of the area as High potential, 40\%
as Medium, and 40\% as Low, with high potential concentrated along the Soan basin
and the northern and riverine alluvial corridors.

\section{Contributions}
\label{sec:concl_contributions}

The study delivers the first machine-learning groundwater potential map of the
full \studyarea{} at 30~m resolution, expanding the feature set of the reference
study from five layers to nine, providing a systematic comparison of four
classifiers on a balanced dataset, and---perhaps most importantly for a
data-scarce region---demonstrating a transparent, evidence-based way to validate
such a map when field measurements are unavailable. The accompanying GeoTIFF
outputs, both a hard classification and a continuous probability surface, are
directly usable in the GIS workflows of regional water-resource agencies.

\section{Limitations}
\label{sec:concl_limitations}

The findings should be read with the limitations set out in
Chapter~\ref{ch:discussion}. Chief among them is label circularity: because the
training labels were derived from the same features used as inputs, the high
internal accuracy measures fidelity to the weighted overlay rather than to ground
truth, and the map's real-world credibility rests on the spatial-realism
validation rather than on its accuracy figures. The validation, though
independent, is built on relatively few points and establishes large-scale
realism rather than pixel-level correctness. The feature set omits subsurface
lithology and lineaments---likely causing over-prediction in hard-rock terrain---
as well as water-quality information, and its static layers describe an average
rather than a time-varying condition.

\section{Recommendations for Future Work}
\label{sec:concl_future}

Several extensions follow naturally and would strengthen the map in roughly this
order of priority.

\begin{enumerate}
  \item \textbf{Add a lithological layer.} Digitizing geological and lineament
        information from Geological Survey of Pakistan sheets and adding it as a
        tenth feature is the single most valuable improvement. It would directly
        address the expected over-prediction in the crystalline and folded terrain
        of the Salt Range, Margalla Hills, and Kala Chitta Range, where structure
        rather than surface condition governs groundwater \cite{benjmel2020morocco}.

  \item \textbf{Incorporate field-measured well data.} Where borehole yield or
        water-level records can be obtained---for example through the PCRWR
        monitoring network---they would allow a true accuracy assessment to
        complement, and eventually move beyond, the knowledge-based labels,
        loosening the circularity that constrains the present interpretation.

  \item \textbf{Add a water-quality dimension.} A contamination or salinity layer
        would let the map distinguish high-potential ground that is fit for use
        from high-potential ground that is not, a distinction that matters
        particularly in the peri-urban Rawalpindi belt \cite{waseem2024waterquality}.

  \item \textbf{Capture temporal dynamics.} Replacing the static rainfall and
        spectral summaries with multi-temporal inputs would allow seasonal and
        interannual variation in recharge to be represented, rather than a single
        average condition.

  \item \textbf{Deliver an interactive product.} Publishing the map as a web-based
        GIS layer would put it directly in the hands of district planners and
        farmers, turning a research output into a usable decision-support tool.
\end{enumerate}

\section{Concluding Remarks}
\label{sec:concl_remarks}

Groundwater will only become more central to the \studyarea{} as rainfall grows
less reliable and demand continues to rise, a pressure visible across aquifers
worldwide \cite{jasechko2024gwdecline}. Knowing where the resource is likely to be
found is a precondition for managing it well. This thesis has shown that machine
learning can take over the mapping task from manual weighted overlay across the
whole plateau, at fine resolution, and---when its limitations are acknowledged
rather than hidden---can be validated against independent hydrogeological evidence
rather than internal accuracy alone. The result is not a final answer but a
defensible and extensible foundation: a first regional, reproducible groundwater
potential map for the \studyarea{}, and a clear path, beginning with the addition
of subsurface geology, toward making it better.
